\documentclass[11pt,a4paper]{article}
\usepackage[italian]{babel}
\usepackage[margin=2.5cm]{geometry}
\usepackage{parskip}
\usepackage{amsmath, amssymb, amsthm}
\usepackage{booktabs}
\usepackage{array}
\usepackage{xcolor}
\usepackage{needspace}
\usepackage{enumitem}
\usepackage{listings}
\usepackage{tikz}
\usepackage[most]{tcolorbox}
\usepackage[hidelinks]{hyperref}
\usetikzlibrary{decorations.pathreplacing, shapes.geometric, arrows.meta, positioning}

% Niente vedove/orfane: i paragrafi non lasciano righe isolate a fine/inizio pagina
\widowpenalty=10000
\clubpenalty=10000

% Elenchi più compatti
\setlist{itemsep=2pt, parsep=0pt, topsep=4pt, partopsep=0pt}

% --- Riquadri con bordo nero, senza numero (si spezzano tra due pagine) ---
% Uso: \begin{definizione} ... \end{definizione}
%      \begin{teorema}[Nome] ... \end{teorema}  (il nome è facoltativo)
\tcbset{nota/.style={enhanced jigsaw, breakable, colback=white,
  colframe=black, boxrule=0.5pt, sharp corners}}
\NewTColorBox{definizione}{}{nota,
  before upper={\textbf{Definizione.}\ }}
\NewTColorBox{teorema}{o}{nota,
  before upper={\textbf{Teorema\IfValueT{#1}{ (#1)}.}\ }}
\NewTColorBox{proposizione}{o}{nota,
  before upper={\textbf{Proposizione\IfValueT{#1}{ (#1)}.}\ }}
\NewTColorBox{proprieta}{o}{nota,
  before upper={\textbf{Proprietà\IfValueT{#1}{ (#1)}.}\ }}
\NewTColorBox{assioma}{o}{nota, boxrule=1pt,
  before upper={\textbf{Assioma\IfValueT{#1}{ (#1)}.}\ }}

% --- Ambienti semplici (senza riquadro) ---
% Il titolo sta in cima al blocco, anche se il contenuto inizia con un riquadro o
% con codice; e non resta mai da solo in fondo a una pagina.
% Uso: \begin{esempio}[Titolo facoltativo] ... \end{esempio}
\newcommand{\newrunin}[2]{%
  \NewDocumentEnvironment{#1}{o}{%
    \par\Needspace{4\baselineskip}\addvspace{\medskipamount}\noindent
    \textbf{#2}\IfValueT{##1}{ (##1)}\textbf{.}\ \ignorespaces}%
  {\par\addvspace{\medskipamount}}}
\newrunin{esempio}{Esempio}
\newrunin{esercizio}{Esercizio}
\newrunin{osservazione}{Osservazione}

% V e F nelle tavole di verità
\newcommand{\V}{\textsf{V}}
\newcommand{\F}{\textsf{F}}

% --- Codice C ---
% Uso: \begin{codice} ... \end{codice}      codice C numerato
%      \begin{comando} ... \end{comando}    comandi da terminale
%      \begin{risultato} ... \end{risultato} output di un programma
%      \lstinline|printf("ciao");|           codice dentro al testo
\lstdefinestyle{c}{
  language=C,
  basicstyle=\ttfamily\small,
  keywordstyle=\color{blue}\bfseries,
  commentstyle=\color{gray}\itshape,
  stringstyle=\color{red!70!black},
  numbers=left,
  numberstyle=\tiny\color{gray},
  numbersep=8pt,
  columns=flexible,
  keepspaces=true,
  breaklines=true,
  showstringspaces=false,
  tabsize=4,
  morekeywords={include, define},
  % lettere accentate nei commenti
  literate={à}{{\`a}}1 {è}{{\`e}}1 {é}{{\'e}}1 {ì}{{\`i}}1
           {ò}{{\`o}}1 {ù}{{\`u}}1 {È}{{\`E}}1 {À}{{\`A}}1
}
\lstset{style=c}
\newtcblisting{codice}{listing only, nota, left=6mm,
  listing options={style=c}}
\newtcblisting{comando}{listing only, nota,
  listing options={basicstyle=\ttfamily\small, numbers=none, language={},
    columns=flexible, keepspaces=true, breaklines=true}}
\newtcblisting{risultato}{listing only, enhanced jigsaw, breakable,
  colback=black!5, colframe=black, boxrule=0.5pt, sharp corners,
  listing options={basicstyle=\ttfamily\small, numbers=none, language={},
    columns=flexible, keepspaces=true, breaklines=true}}

% --- Diagrammi di flusso (simboli standard) ---
\tikzset{
  terminale/.style = {ellipse, draw, minimum width=2.5cm, minimum height=0.9cm},
  io/.style        = {trapezium, trapezium left angle=70, trapezium right angle=110,
                      draw, minimum width=2.5cm, minimum height=0.9cm},
  processo/.style  = {rectangle, draw, minimum width=2.5cm, minimum height=0.9cm},
  decisione/.style = {diamond, draw, aspect=2, minimum width=2.5cm},
  freccia/.style   = {thick, -{Stealth}}
}

\title{Estremi di una funzione e funzioni limitate}
\author{Marco Cerbella}
\date{Lezione 4 -- 5 ottobre 2026}

\begin{document}
\maketitle

\section{Estremo superiore e inferiore di una funzione}

Sia $f : D \to \mathbb{R}$ e $A \subseteq D$.

\begin{definizione}
L'\emph{estremo superiore di $f$ su $A$} è l'estremo superiore dell'insieme $f(A)$:
\[
\sup_{x \in A} f(x) = \sup f(A).
\]
$f$ è \emph{superiormente limitata su $A$} se l'insieme $f(A)$ è superiormente limitato.
\end{definizione}

Se $\sup_{x \in A} f(x)$ è finito e appartiene a $f(A)$, allora è il \emph{valore massimo} di $f$ su $A$. Cioè, se
\begin{enumerate}
    \item per ogni $x \in A$ vale $f(x) \leq M$;
    \item esiste $x_M \in A$ tale che $f(x_M) = M$,
\end{enumerate}
allora $M$ è il valore massimo e $x_M$ è un \emph{punto di massimo}.

\begin{definizione}
L'\emph{estremo inferiore di $f$ su $A$} è l'estremo inferiore dell'insieme $f(A)$:
\[
\inf_{x \in A} f(x) = \inf f(A).
\]
$f$ è \emph{inferiormente limitata su $A$} se $f(A)$ è inferiormente limitato.
\end{definizione}

Se $\inf_{x \in A} f(x)$ è finito e appartiene a $f(A)$, allora è il \emph{valore minimo} di $f$ su $A$. Cioè, se
\begin{enumerate}
    \item per ogni $x \in A$ vale $m \leq f(x)$;
    \item esiste $x_m \in A$ tale che $f(x_m) = m$,
\end{enumerate}
allora $m$ è il valore minimo e $x_m$ è un \emph{punto di minimo}.

\section{Funzioni limitate}

Sia $f : D \to \mathbb{R}$.

\begin{definizione}
\begin{itemize}
    \item $f$ si dice \emph{limitata superiormente} se esiste $M \in \mathbb{R}$ tale che $f(x) \leq M$ per ogni $x \in D$;
    \item $f$ si dice \emph{limitata inferiormente} se esiste $m \in \mathbb{R}$ tale che $f(x) \geq m$ per ogni $x \in D$;
    \item $f$ si dice \emph{limitata} se è limitata sia superiormente sia inferiormente.
\end{itemize}
\end{definizione}

\begin{esempio}
\leavevmode
\begin{enumerate}
    \item $f(x) = x^2$, $D = \mathbb{R}$, $f(\mathbb{R}) = [0, +\infty)$. La funzione è limitata inferiormente ma non superiormente. $m = 0$ è il valore minimo di $f$ e $x_m = 0$ è il punto di minimo.
    \item $f(x) = 3$, $D = \mathbb{R}$, $f(\mathbb{R}) = \{3\}$. La funzione è limitata; $3$ è il valore massimo e il valore minimo, e ogni $x \in \mathbb{R}$ è sia punto di massimo sia punto di minimo.
\end{enumerate}
\end{esempio}

\end{document}
